\documentclass[a4paper,11pt]{article}
\usepackage[utf8]{inputenc}
\usepackage[T1]{fontenc}
\usepackage[french]{babel}
\usepackage[left=2cm,right=2cm,top=2cm,bottom=2cm]{geometry}
\usepackage{xcolor}
\usepackage{hyperref}
\usepackage{fontawesome5}
\usepackage{parskip}

% --- Couleurs Thales ---
\definecolor{primary}{RGB}{0, 45, 114}
\hypersetup{colorlinks=true, urlcolor=primary}

\begin{document}

% --- EXPÉDITEUR ---
\noindent
\begin{minipage}[t]{0.5\textwidth}
    \textbf{Loïc Aron MBASSI EWOLO}\\
    Élève-Ingénieur Bac+5 -- Double diplôme ENSEA\\[3pt]
    \faMapMarker\ Cergy, France\\
    \faPhone\ (+33) 7 59 52 33 98\\
    \faEnvelope\ \href{mailto:loicaron.mbassiewolo@ensea.fr}{loicaron.mbassiewolo@ensea.fr}
\end{minipage}
\hfill
% --- DATE ---
\begin{minipage}[t]{0.4\textwidth}
    \raggedleft
    Cergy, le 11 mars 2026
\end{minipage}

\vspace{1.2cm}

% --- DESTINATAIRE ---
\hfill
\begin{minipage}[t]{0.5\textwidth}
    \raggedleft
    \textbf{Thales}\\
    Digital Services\\
    Élancourt, France
\end{minipage}

\vspace{1cm}

\noindent\textbf{Objet :} Candidature en alternance -- Ingénieur Développement Logiciel Fullstack

\vspace{0.8cm}

Madame, Monsieur,

Trois ans de développement PHP en production, une plateforme de télémédecine maintenue sur deux ans, et un outil open source publié sur Packagist : ce parcours m'a forgé une méthode rigoureuse de conception, de documentation et de livraison. En double diplôme ingénieur à l'ENSEA et l'École Polytechnique de Yaoundé, je suis actuellement en stage de recherche à INRIA Saclay et je serai disponible dès \textbf{septembre 2026} pour rejoindre votre équipe Digital Services à Élancourt en alternance.

Votre besoin de moderniser les outils internes et d'automatiser les processus résonne avec mon expérience. Chez \textbf{CSC} (Fintech), j'ai conçu le backend d'une plateforme de transactions financières en Laravel : APIs RESTful documentées, gestion de concurrence, containerisation Docker et travail en sprints agiles avec démonstrations régulières. Cette mission m'a appris à analyser des besoins métier, proposer des solutions techniques adaptées et livrer de manière itérative, des compétences directement transposables au contexte Thales où vous collaborez avec les services Qualité, HSE, Production et Support.

Sur la plateforme \textbf{SOSDOCTA} (télémédecine, 2 ans de maintenance), j'ai développé l'intégralité du stack : backend PHP/MySQL, dashboards responsifs en JavaScript, système de réservation et intégration de paiements sécurisés. La gestion de données médicales sensibles m'a familiarisé avec les exigences de conformité RGPD et de sécurisation des accès (OAuth2/JWT). La maintenance évolutive (corrections, nouvelles fonctionnalités, montées de version) correspond aux missions de suivi et d'amélioration continue décrites dans votre offre.

Par ailleurs, mon projet \textbf{GoFind} m'a permis de pratiquer l'architecture microservices (Laravel, NestJS, RabbitMQ, Docker) avec intégration continue, et le projet \textbf{Freengo} m'a donné une expérience concrète en Java Spring Boot. Mon \textbf{Laravel API Generator}, outil open source de génération automatique de code backend, témoigne de ma capacité à mener des initiatives R\&D et à concevoir des Proof of Concept de manière autonome.

Je suis disponible pour un entretien afin de vous présenter mon parcours et discuter de ma contribution à vos projets.

Je vous prie d'agréer, Madame, Monsieur, l'expression de mes salutations distinguées.

\vspace{0.8cm}

\textbf{Loïc Aron MBASSI EWOLO}\\
\textit{Élève-Ingénieur ENSEA / Polytechnique Yaoundé}

\end{document}
